% Môn: Toán
% Lớp: 10
% Chương: Chương I. Mệnh đề và tập hợp
% Bài:  Mệnh đề · Xác định mệnh đề và tính đúng sai
% Số câu: 0
% App đọc file này trực tiếp — sửa rồi Commit trên GitHub
% Ghi nguồn từng câu: \nguon{SGK} hoặc \nguon{2 SBT VL 12 KNTT} trong


%%%=============EX_1=============%%%


%%%=============EX_2=============%%%

\dangbt{Xác định mệnh đề và tính đúng sai}
\begin{ex}
	% ID: T10C1B1-02-TN
	% Mức: NB
	Câu nào sau đây là mệnh đề?
	
	\choice
		{Cháu chào ông bà ạ!}
		{Trời hôm nay đẹp quá!}
		{\True Năm $2022$ là năm nhuận.}
		{Chiều nay đi đánh cầu lông không?}
	\loigiai{
		Câu \lq\lq Trời hôm nay đẹp quá!\rq\rqvà câu \lq\lq Cháu chào ông bà ạ! \rq\rqlà câu cảm thán nên không phải là mệnh đề.\\
		Câu \lq\lq Chiều nay đi đánh cầu lông không?\rq\rqlà câu hỏi nên cũng không phải là mệnh đề. \\
		Câu \lq\lq Năm $2022$ là năm nhuận \rq\rqlà câu khẳng định nên là mệnh đề.}
\end{ex}
% Mức: NB
% ID: T10C1B1-04-TN
% ID: T10C1B1-04-TN
% Mức: NB
% ID: T10C1B1-03-TN
% ID: T10C1B1-03-TN
% Mức: NB
% ID: T10C1B1-03-TN

% Mức: NB
% ID: T10C1B1-05-TN
% ID: T10C1B1-05-TN
% Mức: NB
% ID: T10C1B1-04-TN
% ID: T10C1B1-04-TN
% Mức: NB
% ID: T10C1B1-04-TN

% Mức: TH
% ID: T10C1B1-06-TN
% ID: T10C1B1-06-TN
% Mức: TH
% ID: T10C1B1-05-TN
% ID: T10C1B1-05-TN
% Mức: TH
% ID: T10C1B1-05-TN
% Mức: TH

%%%=============EX_3=============%%%
\begin{ex}
	% ID: T10C1B1-05-TN
	% Mức: TH
	Trong các câu sau có bao nhiêu câu là mệnh đề
	\begin{enumEX}{2}
		\item Ngôi nhà đẹp quá!
		\item Năm 1900 không phải là năm nhuận.
		\item Số 9 là số nguyên tố.
		\item $x^3+1=0$.
	\end{enumEX}
	\choice
		{$1$}
		{$3$}
		{\True $2$}
		{$4$}
	\loigiai{
		\lq\lq Ngôi nhà đẹp quá!\rq\rqlà câu cảm thán, không phải là mệnh đề.\\
		\lq\lq $x^3+1=0$\rq\rqlà mệnh đề chứa biến, không phải là mệnh đề.}
\end{ex}
% Mức: TH
% ID: T10C1B1-61-TN
% ID: T10C1B1-61-TN
% Mức: TH
% ID: T10C1B1-60-TN
% ID: T10C1B1-60-TN
% Mức: TH
% ID: T10C1B1-60-TN
% Mức: TH

%%%=============EX_5=============%%%
\begin{ex}
	% ID: T10C1B1-60-TN
	% Mức: TH
	Cho mệnh đề chứa biến $P(x)\colon \lq\lq x^2{<}3x$ với $x$ là số thực. Mệnh đề nào sau đây đúng?
	\choice
		{$P(0)$}
		{$P(5)$}
		{$P(-1)$}
		{\True $P(2)$}
	\loigiai{
		Thay các giá trị $2, 0, 5, -1$ ta được $2^2{<}3\cdot 2$ nên $P(2)$ đúng.
		}
\end{ex}
% Mức: NB
% ID: T10C1B1-65-TN
% ID: T10C1B1-65-TN
% Mức: NB
% ID: T10C1B1-64-TN
% ID: T10C1B1-64-TN
% Mức: NB
% ID: T10C1B1-64-TN

% Mức: NB
% ID: T10C1B1-66-TN
% ID: T10C1B1-66-TN
% Mức: NB
% ID: T10C1B1-65-TN
% ID: T10C1B1-65-TN
% Mức: NB
% ID: T10C1B1-65-TN

% Mức: NB
% ID: T10C1B1-67-TN
% ID: T10C1B1-67-TN
% Mức: NB
% ID: T10C1B1-66-TN
% ID: T10C1B1-66-TN
% Mức: NB
% ID: T10C1B1-66-TN

% Mức: NB
% ID: T10C1B1-68-TN
% ID: T10C1B1-68-TN
% Mức: NB
% ID: T10C1B1-67-TN
% ID: T10C1B1-67-TN
% Mức: NB
% ID: T10C1B1-67-TN

% Mức: TH
% ID: T10C1B1-69-TN
% ID: T10C1B1-69-TN
% Mức: TH
% ID: T10C1B1-68-TN
% ID: T10C1B1-68-TN
% Mức: TH
% ID: T10C1B1-68-TN

% Mức: NB
% ID: T10C1B1-70-TN
% ID: T10C1B1-70-TN
% Mức: NB
% ID: T10C1B1-69-TN
% ID: T10C1B1-69-TN
% Mức: NB
% ID: T10C1B1-69-TN

% Mức: TH
% ID: T10C1B1-71-TN
% ID: T10C1B1-71-TN
% Mức: TH
% ID: T10C1B1-70-TN
% ID: T10C1B1-70-TN
% Mức: TH
% ID: T10C1B1-91-TN
% Mức: NB

%%%=============EX_8=============%%%
\begin{ex}
	% ID: T10C1B1-91-TN
	% Mức: NB
	Mệnh đề là:
	\choice
		{Một khẳng định luôn đúng}
		{Câu cảm thán}
		{Câu nghi vấn hoặc câu cầu khiến}
		{\True Một khẳng định chỉ có thể đúng hoặc sai}
	\loigiai{
		Mệnh đề là một khẳng định chỉ có thể đúng hoặc sai.
		}
\end{ex}

%%%=============EX_9=============%%%
\begin{ex}
	% ID: T10C1B1-216-DS
	% Mức: TH
	Xét tính đúng sai của các khẳng định sau.
	\choiceTF
		{Phát biểu "Ở đây đẹp quá!" là một mệnh đề}
		{\True Phát biểu "Phương trình $x^2-3 x+1=0$ vô nghiệm." là một mệnh đề}
		{\True Phát biểu "Số $16$ không phải là số nguyên tố." là một mệnh đề}
		{Phát biểu "Số $\pi$ có lớn hơn $3$ hay không?" là một mệnh đề}
	\loigiai{
		\begin{itemchoice}
			\itemch sai.
			\itemch đúng.
			\itemch đúng.
			\itemch đúng.
		\end{itemchoice}
		}
\end{ex}
%%%=============EX_10=============%%%
\begin{ex}
	% ID: T10C1B1-217-TLN
	% Mức: TH
	Trong các mệnh đề sau, có bao nhiêu mệnh đề toán học?  $\bullet$ $A\colon$ '' Hai đường thẳng song song không có điểm chung ''.  $\bullet$ $B\colon$ '' Tích hai số thực trái dấu là một số thực âm ''.  $\bullet$ $C\colon$ '' Hàm số $y = x^2$ có đồ thị là một parabol ''.  $\bullet$ $D\colon$ '' Mặt Trời quay quanh Trái Đất ''.  $\bullet$ $E\colon$ '' Mọi số tự nhiên đều là dương ''.
	\shortans{4}
\loigiai{$\bullet$ Phát biểu '' Hai đường thẳng song song không có điểm chung '' là một mệnh đề toán học.  $\bullet$ Phát biểu '' Tích hai số thực trái dấu là một số thực âm '' là một mệnh đề toán học.  $\bullet$ Phát biểu '' Hàm số $y = x^2$ có đồ thị là một parabol '' là một mệnh đề toán học.  $\bullet$ Phát biểu '' Mặt Trời quay quanh Trái Đất '' không là một mệnh đề toán học (vì không liên quan đến sự kiện Toán học nào).  $\bullet$ Phát biểu '' Mọi số tự nhiên đều là dương '' là một mệnh đề toán học.}
\end{ex}
%%%=============EX_11=============%%%
\begin{ex}
	% ID: T10C1B1-105-TLN
	% Mức: TH
	Có bao nhiêu câu sau đây là mệnh đề ?  $\bullet$ $A\colon$ '' Ôi! Trời lạnh thật đấy! ''.  $\bullet$ $B\colon$ '' Trái Đất hình lập phương ''.  $\bullet$ $C\colon$ '' Hàm số $y = x^2$ có đồ thị là một parabol ''.  $\bullet$ $D\colon$ '' Pi là một số vô tỷ ''.  $\bullet$ $E\colon$ '' Mọi số tự nhiên đều là dương ''.
	\shortans{4}
\loigiai{$\bullet$ Phát biểu '' Ôi! Trời lạnh thật đấy! '' không phải là một mệnh đề vì không thể xác định đúng hoặc sai.  $\bullet$ Phát biểu '' Trái Đất hình lập phương '' là một mệnh đề vì có thể xác định đúng hoặc sai.  $\bullet$ Phát biểu '' Hàm số $y = x^2$ có đồ thị là một parabol '' là một mệnh đề vì có thể xác định đúng hoặc sai.  $\bullet$ Phát biểu '' Pi là một số vô tỷ '' là một mệnh đề vì có thể xác định đúng hoặc sai.  $\bullet$ Phát biểu '' Mọi số tự nhiên đều là dương '' là một mệnh đề vì có thể xác định đúng hoặc sai.}
\end{ex}